\documentclass{article}
\usepackage{arxiv}

\usepackage[utf8]{inputenc}
\usepackage[T1]{fontenc}
\usepackage{booktabs}
\usepackage{amsfonts}
\usepackage{nicefrac}
\usepackage{microtype}
\usepackage{graphicx}
\usepackage{natbib}
\usepackage{doi}

% Generated numerical constants from results/numbers.json
% Will be populated during experiment execution at Gate 5.
\newcommand{\DummyPlaceholder}{TBD}


\title{Does the Added Prognostic Value of RNA-seq Over Clinical Variables Survive Site-Held-Out Validation? A Pre-Specified Re-Evaluation in TCGA-KIRC}

\author{
  Vikhram S \\
  Department of Biotechnology \\
  Saveetha Institute of Medical and Technical Sciences \\
  Chennai, India \\
  \texttt{vikhrams@saveetha.ac.in}
}

\date{October 1, 2026}

\begin{document}
\maketitle

\begin{abstract}
\textbf{Background:} Prognostic multivariable prediction models combining bulk transcriptomic data with baseline clinical predictors frequently report substantial gains in discrimination under randomized cross-validation. However, multi-center observational datasets may exhibit center-level acquisition variability and batch confounding.
\textbf{Methods:} In this pre-specified study in TCGA Kidney Renal Clear Cell Carcinoma (TCGA-KIRC), we evaluate whether incremental prognostic value ($\Delta C$) survives site-held-out validation. We fit an unpenalized clinical Cox model (M0: age, sex, stage, grade) and an incremental model (M1: unpenalized clinical covariates plus elastic-net penalized top-variance RNA-seq features). Discrimination is evaluated across repeated stratified random cross-validation versus repeated site-grouped cross-validation partitioned by Tissue Source Site (TSS).
\textbf{Results:} Results pending experimental execution following protocol freeze.
\textbf{Conclusions:} Study conclusions will reflect empirical findings exactly as observed without selective reporting.
\end{abstract}

\keywords{Survival Analysis \and Prognostic Biomarkers \and TCGA-KIRC \and Site-Held-Out Validation \and Batch Effects}

\section{Introduction}
Prognostic prediction models utilizing high-throughput transcriptomics have proliferated across oncology literature. An essential question in translational biomarker development is whether molecular profiles provide incremental discrimination over readily available clinical and pathological covariates.

\section{Related Work}
Prior investigations have documented significant site-associated variations and center-correlated artifacts in multi-center datasets \citep{howard2021site, herrmann2020benchmark, survboard2023}.

\section{Methods}
\subsection{Cohort and Study Design}
The study cohort comprises primary solid tumor samples from TCGA-KIRC obtained via UCSC Xena.

\subsection{Cross-Validation Regimes}
We compare repeated event-stratified cross-validation against repeated site-grouped cross-validation using Tissue Source Site identifiers.

\subsection{Model Specifications}
M0 denotes the clinical baseline Cox proportional hazards model. M1 combines clinical covariates and top-variance RNA features regularized via elastic-net penalties.

\subsection{Statistical Analysis}
Primary endpoint is paired $\Delta C = C_{\text{M1}} - C_{\text{M0}}$. Uncertainty is characterized via clustered bootstrap resampling.

\section{Results}
\textit{[Results will be generated programmatically from experimental outputs.]}

\section{Discussion}
\textit{[Discussion will be formulated based strictly on empirical results.]}

\section{Limitations}
Key limitations include retrospective observational registry constraints, proxy site definitions, and non-interventional design.

\section{Declarations}
\subsection*{Ethics Approval}
The study uses publicly available, de-identified retrospective data from TCGA; formal ethical review waiver applies [wording TO VERIFY with medRxiv].

\subsection*{Data and Code Availability}
Code is available at \url{https://github.com/Vikhram-S/kirc-site-heldout-survival}. Raw data access procedures are documented via UCSC Xena. [Zenodo DOI TO FILL].

\subsection*{Competing Interests}
The author declares no competing interests.

\subsection*{Funding}
No external funding was received for this work.

\subsection*{Contributions}
V.S. conceived the study, specified the protocol, executed analysis, and authored the manuscript.

\subsection*{AI-Assistance Disclosure}
LLM-based coding assistance was used for code scaffolding and document formatting. All scientific design and methodological decisions were directed by the human author.

\bibliographystyle{unsrtnat}
\bibliography{references}

\end{document}
